\section{From Assigned Roles to Realized Organization}
\label{sec:framework}

\subsection{Specification and realization}

An organization is specified by a tuple: a set of agents, an interaction order,
a message-passing graph, a context-visibility policy, and -- the part of
interest here -- a \emph{role contract} per agent. A role contract is a set of
clauses, each of which asserts either that the agent performs some function
(\textsc{must} decompose the problem; \textsc{must} state the required answer
form) or that it does not (\textsc{must not} carry out arithmetic;
\textsc{must not} state a final answer). Clauses of the second kind are what
make the specification a \emph{division} of labor rather than a list of
suggestions: they assign work away from one agent and, implicitly, to another.

The realized organization is the behavior that occurs when the specification is
executed by a particular actor substrate. Because a contract clause is prose
handed to a language model, realization is not guaranteed at any clause, and
nothing in the specification records which clauses were honored. An evaluation
that observes only the task outcome cannot distinguish an organization that
divided the work as specified from one in which a single agent did everything
while the others produced well-formed commentary. Both appear as a score.

\subsection{Two constructs that must not be merged}

\paragraph{Role distinguishability.} Can an observer recover which role emitted
a message from the message's behavior alone? This is a property of the induced
behavioral distributions: it asks whether the roles differ, not whether they
differ \emph{in the way specified}. It is cheap to measure, it is what a reader
usually has in mind when asking whether role prompts ``work,'' and -- as
\S\ref{sec:measurement} shows -- it can be very high while the specified
division is almost entirely absent.

\paragraph{Role fidelity.} Does an agent perform the function its own contract
assigns? This is defined only relative to a contract, it is evaluated
clause-by-clause, and it is not comparable across agents that received
different contracts. Fidelity is the construct that carries the organizational
content, because it is the one that references the specification.

The two are logically independent. An agent that ignores every clause of its
contract but writes in a recognizable house style is perfectly distinguishable
and has zero fidelity. An agent that honors every clause while formatting its
output exactly like its neighbor has full fidelity and is indistinguishable.
Conflating them makes it possible to report that roles ``are working'' on the
basis of a measurement that is insensitive to whether the division of labor
occurred, and \S\ref{sec:measurement} shows that this is not a hypothetical
concern in the systems we study.

\subsection{Actor-dependence and the resulting design}

Because realization occurs at the actor, the same specification executed by
different models is potentially a different organization. This has a
methodological consequence that is easy to state and easy to violate: a claim
about an organization that is measured under one actor is a claim about that
actor--organization pair. It also has a design consequence, which is the
remainder of this paper. If a specification is advice, then the intervention
that tests its role is not a more emphatic specification but a structural one
-- an interface under which the violation cannot be expressed. Whether such an
intervention changes the realized organization, whether it changes
effectiveness, and whether either change is actor-dependent are the questions
\S\ref{sec:method} operationalizes.
